
\section{\trCM{}环与概形}

\textbf{26.2.1.}\ \emph{\trhDefinition。}\quad \trNoether{}局部环 $(A,\mathfrak{m})$ 称为\textbf{\trCM}的，若 $\operatorname{depth}A=\dim A$，即在 \eqref{eq:depth-le-dim} 中等号成立。（类似地也可以定义\trCM{}模，但我们用不到这个概念。）一个\trLocally\trNoether{}\trScheme{}称为\textbf{\trCM}的，若它的所有\trLocalRing{}都是\trCM{}的。我们不会使用这种说法，不过：环 $A$ 称为\textbf{\trCM{}环}，若它在所有素理想 $\mathfrak{p}$ 处的局部化 $A_{\mathfrak{p}}$ 都是\trCM{}的——等价地，若 $\Spec A$ 是\trCM{}的。（下面的\mbox{事实~26.2.14} 说：局部环的\trCM{}性在局部化下保持。\footnote{原书此处作「Cohen-Maculayness」，脱一字母「a」，当为「Cohen-Macaulayness」之误，径改。——译者注}这便保证了此定义与\trCM{}局部环的定义之间不存在矛盾。）

\textbf{26.2.A.（练习）}\quad 证明：每个维数为 $0$ 的\trLocally\trNoether{}\trScheme{}都是\trCM{}的。再证明：维数为 $1$ 的\trLocally\trNoether{}\trScheme{}是\trCM{}的，当且仅当它没有\trEmbeddedPoint。

\textbf{26.2.2.}\ \emph{（反）例。}\quad 设 $X$ 是\mbox{练习~26.1.G} 的那个例子——交于一点的两个平面。由\mbox{练习~26.1.G}，$X$ 不是\trCM{}的。

\textbf{26.2.B.（练习）}\quad 设 $A$ 是\trCM{}的\trNoether{}局部环。用\mbox{定理~26.1.7} 证明：$\Spec A$ 是\trPureDimensional{}的，并且没有\trEmbeddedPoint。（\trNoether{}局部环即使有\trPureDimension{}且无\trEmbeddedPrime，也未必是\trCM{}的；见\mbox{（反）例~26.2.2}。）

\textbf{26.2.3. \trhTheorem（\trCM{}性的\trSlicing{}判别法）。}\ ——\ \emph{设 $(A,\mathfrak{m})$ 是\trNoether{}局部环，$x\in\mathfrak{m}$ 是一个\trNonZerodivisor。则 $(A,\mathfrak{m})$ 是\trCM{}的，当且仅当 $(A/x,\mathfrak{m})$ 是\trCM{}的。}

把它与正则性及平坦性的\trSlicing{}判别法相比较（分别是\mbox{练习~13.2.C} 与\mbox{定理~24.6.6}）。

\textbf{26.2.4.}\ \emph{\trSlicing{}判别法的几何解释。}\quad 设 $X$ 是\trLocally\trNoether{}\trScheme，$D$ 是一个\trEffectiveCartierDivisor。若 $X$ 是\trCM{}的，则 $D$ 亦然。若 $D$ 是\trCM{}的，则 $X$ 在 $D$ 的各点处都是\trCM{}的。

\textbf{26.2.C.（练习）}\quad 证明\mbox{定理~26.2.3}：用\mbox{定理~26.1.3}，以及「极大\trRegularSequence（在 $\mathfrak{m}$ 中的）都具有相同长度」这一事实。（提示：回想\mbox{练习~12.3.I}：对任意 $f\in A$ 都有 $\dim A/(f)\ge\dim A-1$。）

\textbf{26.2.D.（练习）}\quad 证明：若 $(A,\mathfrak{m})$ 是\trCM{}的，则一组元素 $x_1,\dots,x_r\in\mathfrak{m}$ 构成（$A$ 的）\trRegularSequence，当且仅当 $\dim A/(x_1,\dots,x_r)=\dim A-r$。

\textbf{26.2.5.}\quad 由\mbox{练习~26.1.F}，\trRegular\trLocalRing（按定义是诺特的）都是\trCM{}的。特别地，由于域 $k$ 上的\trSmooth{}概形是正则的（\mbox{定理~13.2.7(b)}，其证明见 \S24.8.9），我们看到光滑 $k$-\trScheme{}都是\trCM{}的。（结合\mbox{练习~26.2.B}，这便解释了光滑\trVariety{}上的\trEffectiveCartierDivisor{}为何没有\trEmbeddedPoint，印证了\mbox{题外话~20.2.1} 中的一处评注。）将此事与\mbox{练习~26.2.D}（或 \S26.2.4）相结合，便得到下述结论。

\textbf{26.2.6. \trhProposition。}\ ——\ \emph{\trSmooth{} $k$-\trScheme{}中的\trRegularEmbedding{}都是\trCM{}的。}

\textbf{26.2.7.}\quad 作为\mbox{命题~26.2.6} 与「\trCM{}概形没有\trEmbeddedPoint」（\mbox{练习~26.2.B}）的推论，我们看到：\trSmooth{} $k$-\trScheme{}（比如 $\mathbb{A}^n_k$ 或 $\mathbb{P}^n_k$）中的\trRegularEmbedding{}都没有\trEmbeddedPoint，这推广了\mbox{练习~6.6.F}（$\mathbb{A}^n_k$ 中超曲面的情形）。若没有\trCM{}性理论，这一点并不显然！

\textbf{26.2.8.}\ \emph{\trCM{}性的另一种定义。}\quad \trSlicing{}判别法（\mbox{定理~26.2.3}）给出了一种富有启发性的\trCM{}性归纳定义，改用\trEffectiveCartierDivisor{}来表述，其精神与 \S26.1.5 计算\trDepth{}的方法如出一辙。暂且设 $(A,\mathfrak{m})$ 是\trNoether{}局部环。
\begin{itemize}
\item[（i）] 若 $\dim A=0$，则 $A$ 是\trCM{}的（\mbox{练习~26.2.A}）。
\item[（iia）] 若 $\dim A>0$，而 $\mathfrak{m}$ 的每个元素都是\trZerodivisor，则 $A$ 不是\trCM{}的（由\mbox{练习~26.1.A}）。
\item[（iib）] 否则，在 $\mathfrak{m}$ 中任取一个\trNonZerodivisor{} $x$。则 $A$ 是\trCM{}的，当且仅当 $A/(x)$（由\mbox{\trKrull{}\trPrincipalIdealTheorem~12.3.3}，其维数必为 $\dim A-1$）是\trCM{}的。
\end{itemize}

下面这个例子，其实早在我们知道任何代数几何之前就可以陈述出来（只是无法证明）。（我们取 $\mathbb{C}$ 而非任意的域来工作，只是为了让这个命题所需的背景知识尽可能少。）

\textbf{26.2.E.（趣味练习：马克斯·诺特（Max Noether）的 AF$+$BG 定理）}\quad 设 $f,g\in\mathbb{C}[x_0,x_1,x_2]$ 是两个\trHomogeneous{}多项式，它们切出 $\mathbb{P}^2_{\mathbb{C}}$ 中的两条曲线，并「横截地」相交，即交于有限多个既约点。设 $h\in\mathbb{C}[x_0,x_1,x_2]$ 是在这些点处取零值的\trHomogeneous{}多项式。证明：$h\in(f,g)$。提示：证明 $\mathbb{A}^3$ 中仿射锥 $V(f)$ 与 $V(g)$ 的交没有\trEmbeddedPoint。（不借助本章发展的理论，这道题做起来相当不平凡！为示此事之精微：你很容易构造三个二次型 $e,f,g\in\mathbb{C}[x_0,x_1,x_2]$，使它们恰好切出 $\mathbb{P}^2_{\mathbb{C}}$ 中的两个点 $[1,0,0]$ 与 $[0,1,0]$，然而直线 $z=0$ 却因次数原因不在理想 $(e,f,g)$ 之中。）

作为\trCM{}性的又一重动机：下一道练习看似基础，但要把证明做得干净，我们需要\trCM{}性。

\textbf{26.2.F.（练习）}\quad 设 $H_1,\dots,H_n$ 是 $\mathbb{P}^n_{\mathbb{C}}$ 中次数分别为 $d_1,\dots,d_n$ 的 $n$ 张\trHypersurface，并设它们的\trSchemeTheoretic{}交 $Z$ 是有限个点。证明：$Z$ 在诸点处的长度之和（对 $Z$ 的诸点求和）为 $d_1\cdots d_n$。等价地：证明 $Z$ 的\trLength{}为 $d_1\cdots d_n$。不严格地说：这些\trHypersurface{}相交于 $d_1\cdots d_n$ 个点，「按正确方式计数」。提示：\trBezout{}定理，取\mbox{练习~18.6.J} 的形式。

\textbf{26.2.9. \trMiracleFlat。}

我们以一条惊人地简洁而有用的平坦性判别法来收束本节，我们将在\trSerreDuality{}的证明中用到它。其核心内容是下述代数结果。

\textbf{26.2.10. \trhTheorem（\trMiracleFlat{}定理，代数版本）。}\ ——\ \emph{设 $\phi\colon(B,\mathfrak{n})\to(A,\mathfrak{m})$ 是\trNoether{}局部环之间的一个（局部）同态，使得 $A$ 是\trCM{}的，$B$ 是正则的，而 $A/\mathfrak{n}A=A\otimes_B(B/\mathfrak{n})$（对应于该纤维的环）具有纯维数 $\dim A-\dim B$。则 $\phi$ 是平坦的。}

\emph{证明。}\quad 我们对 $\dim B$ 作归纳来证明\mbox{定理~26.2.10}。若 $\dim B=0$，则 $B$ 是一个域（\mbox{练习~13.2.A}），于是结论立得，因为域上的任何东西都平坦（\mbox{练习~24.1.A(a)}）。再设 $\dim B>0$，并假设结论已对一切「维数更小的 $B$」得证。取 $x\in\mathfrak{n}\setminus\mathfrak{n}^2$，于是 $B/x$ 是维数为 $\dim B-1$ 的\trRegular\trLocalRing（\mbox{练习~12.3.I}）。那么
\begin{align*}
\dim A/xA &\le \dim B/(x)+\dim A/\mathfrak{n}A &&\text{（关键练习~12.4.A）}\\
&= \dim B-1+\dim A/\mathfrak{n}A &&\text{（练习~12.3.I）}\\
&= \dim A-1 &&\text{（由定理~26.2.10 的假设）}.
\end{align*}
由（\mbox{练习~12.3.I}）可知 $\dim A/xA\ge\dim A-1$，于是我们有 $\dim A/xA=\dim B/(x)+\dim A/\mathfrak{n}A=\dim A-1$。由\mbox{练习~26.2.D}，$A/xA$ 是\trCM{}环，而 $x$ 在 $A$ 上不是\trNonZerodivisor。归纳假设随之适用于 $(B/(x),\mathfrak{n})\to(A/xA,\mathfrak{m})$，故 $A/xA$ 在 $B/(x)$ 上平坦。再由平坦性的局部\trSlicing{}判别法（在靶上，\mbox{定理~24.6.6}），$B\to A$ 是平坦的，此即所欲证。\hfill$\square$

\textbf{26.2.11. \trhTheorem（\trMiracleFlat{}定理）。}\ ——\ \emph{设 $\pi\colon X\to Y$ 是\trPureDimensional{}的\trFiniteType{} $k$-\trScheme{}之间的一个\trMorphism，其中 $X$ 是\trCM{}的，$Y$ 是正则的，而 $\pi$ 的诸\trFiber{}都具有维数 $\dim X-\dim Y$。则 $\pi$ 是平坦的。}

\textbf{26.2.G.（练习）}\quad 证明\trMiracleFlat\mbox{定理~26.2.11}。（别忘了概形通常还有非闭点！）

\trMiracleFlat\mbox{定理~26.2.11} 中的几何局面，是下面这个漂亮「套装」（package）的一部分。

\textbf{26.2.H.（练习：\trMiracleFlat{}套装）}\quad 设 $\pi\colon X\to Y$ 是\trLocally\trFiniteType{} $k$-\trScheme{}之间的一个映射，其中 $X$ 与 $Y$ 都是\trPureDimensional{}的，且 $Y$ 是正则的。证明：下列三条中若任何两条成立，则第三条也成立：
\begin{itemize}
\item[（i）] $\pi$ 是相对维数 $\dim X-\dim Y$ 的平坦\trMorphism。
\item[（ii）] $X$ 是\trCM{}的。
\item[（iii）] 每个\trFiber{} $X_y$ 都是\trCM{}的，且具有纯维数 $\dim X-\dim Y$。
\end{itemize}
提示：若 $\phi\colon B\to A$ 是平坦的环映射，则 $\phi$ 把\trNonZerodivisor{}送到\trNonZerodivisor（\mbox{观察~24.1.2}）。

\mbox{练习~26.2.H} 的陈述还可以改进，代价是牺牲对称性。在「（ii）与（iii）蕴含（i）」这一方向中，（iii）里的\trCM{}假设并不需要。

\textbf{26.2.12.}\ \emph{\trhExample。}\quad 作为\mbox{练习~26.2.H} 实战的一个例子，我们考虑交于一点的两个平面这一例子，沿用\mbox{练习~26.1.G} 的记号。\footnote{原书此处作「two planes meeting a point」，脱「at」，当为「meeting at a point」之误，径补。——译者注}考虑由 $(w-y,\,x-z)$ 给出的\trMorphism{} $\pi\colon X\to Y:=\mathbb{A}^2$。注意 $Y$ 是正则的（\S13.3.6）。但 $X$ 不是\trCM{}的（\mbox{练习~26.1.G}），且 $\pi$ 不是平坦的（\mbox{练习~24.3.I}），于是由\mbox{练习~26.2.H}，你可以用其中任意一条去证明另一条。

\textbf{26.2.13. \trCM{}性的花哨性质。}

我们再列举\trCM{}性的几条附加性质，不加证明。它们值得一看，但我们不会用到。

设 $A$ 是诺特环，且是某个\trCM{}环的商环，则 $\Spec A$ 中\trCM{}点构成的\trLocus{}是开的；见 [Mat2, Ex. 24.2]。与此相关：

\textbf{26.2.14. 事实（参看\mbox{事实~13.8.5}）。}\ ——\ \emph{\trCM{}\trNoether{}局部环在任一\trPrimeIdeal{}处的局部化仍是\trCM{}的。}

（直接证明见 [Stacks, tag 00NB]、[E, Prop. 18.8]、[Mat2, Thm. 17.3(iii)] 或 [Mat1, Thm. 30]。）特别地，对诸簇而言，\trCM{}性可以在闭点处检验（\mbox{练习~5.3.F}）。利用「\trCM{}性在局部化下保持」这一事实，还可以很快证明：\trCM{}局部环都是\trCatenary{}的（见 [E, Cor. 18.10]、[Mat1, Thm. 31(ii)] 或 [Mat2, Thm. 17.9]）。环 $A$ 是\trCM{}的，当且仅当 $A[x]$ 是\trCM{}的（[E, Prop 18.9]、[Mat2, Thm 17.7]、[Mat1, Thm. 33]），当且仅当 $A[[x]]$ 是\trCM{}的 [Mat2, p. 137]。局部环是\trCM{}的，当且仅当它的\trCompletion{}是\trCM{}的（[E, Prop. 18.8]、[Mat2, Thm 17.5]）。

至于这些以及许多其他令人愉快的事实，还有一般地对\trCM{}性的一份透彻入门，见 [BruH]。
